\documentclass[10pt,a4paper]{amsart}
\usepackage[margin=19mm]{geometry}
\usepackage{amsmath,amssymb,mathtools}
\usepackage[scr=rsfs,cal=euler]{mathalfa}
\usepackage{xfrac}
\usepackage[unicode,hidelinks,bookmarks=false]{hyperref}
\newcommand{\ie}{i.e.\ }
\newcommand{\cf}{cf.\ }
\newcommand{\resp}{respectively\ }
\newcommand{\Subsection}{\subsection}
\providecommand{\nobreakdash}{\nobreak\hbox{-}}
\setlength{\emergencystretch}{2em}
\multlinegap0pt
\usepackage[shortlabels]{enumitem}
\usepackage{amssymb}
\newtheorem{theorem}{Theorem}
\newcommand{\Ish}{\mathcal{I}}
\newcommand{\K}{\mathrm{K}}      
\newcommand{\QQ}{\mathbb{Q}} % rational numbers
\newcommand{\ZZ}{\mathbb{Z}} % integers
\renewcommand{\leq}{\leqslant}
\title{\uppercase{CM}-line bundles and slope $\K$-semistability for big and nef line bundles along subschemes}
\author{Nathan Grieve}
\date{Source: arXiv:2509.17546v2 (CC BY 4.0)}
\begin{document}
\maketitle

\noindent\emph{Source and licence.} \uppercase{CM}-line bundles and slope $\K$-semistability for big and nef line bundles along subschemes. Version arXiv:2509.17546v2 (CC BY 4.0), distributed by arXiv under the Creative Commons Attribution 4.0 International licence, http://creativecommons.org/licenses/by/4.0/, as recorded in the arXiv licence metadata for this exact version and shown on the abstract page https://arxiv.org/abs/2509.17546v2. Full licence notice: \texttt{licenses/arxiv-2509.17546-CC-BY-4.0.txt}.\par\smallskip

\noindent\textit{Editorial scope.} Counted unit: the article's theorem Ross:Thomas:Slope:Stab:DF (source0.tex lines 323--342). The article's complete original proof of it is delivered with it (source0.tex lines 344--375, one \texttt{proof} environment of 32 lines). No mathematics is written, repaired or re-derived: the statement and the proof are the article's own lines, in the article's own notation, and no retained line is substituted or re-flowed.  Retained uncounted as the source-local context the proof needs: lines 237--240; lines 260--262.  Nothing was omitted from any retained span, so the package carries no omission record.  The retained lines cite 3 works, whose entries are transcribed verbatim from the article's own bibliography source in the e-print and delivered with short editorial labels recorded in the item's reference mapping.  No retained line contains a figure environment, an image-inclusion command or drawing code, so no image is delivered and the package declares no asset dependency.  Editorial formatting only, in addition: an AMS \texttt{amsart} preamble that restates the article's own declarations and macros used by the retained lines, namely the environments \texttt{theorem} on separate counters, together with the standard packages those lines need.  The retained environments print in this document as Theorem 1 in order of appearance, whereas the article prints the same items with the numbering of its own full document.  The complete native source of the article as distributed in the arXiv e-print for this exact version is bundled unchanged as \texttt{sources/185462/source0.tex}.
\begin{theorem}\label{CM:line:bundle:epsilon:zero}
Let $f \colon X \rightarrow B$ be a flat family of $n$-dimensional projective schemes.  Suppose that $L$ is a big and nef line bundle on $X$.  Fix an ample line bundle $A$ on $X$.  Then
$$\lambda_{\mathrm{CM}}(X,L) = \lim\limits_{\epsilon \to 0} \lambda_{\mathrm{CM}}(X,L \otimes A^{\otimes \epsilon})  \text{.} $$
\end{theorem}

\begin{theorem}[Paul and Tian; Arezzo-et-al {\cite[Proposition 3.3]{Arezzo:DellaVedova:LaNave}}]\label{Paul:Tian:DF:Thm}
Let $(\mathcal{X},\mathcal{L})$ be a big and nef test configuration for $L$ a big and nef line bundle on a projective variety $X$.  Then the Donaldson-Futaki invariant $\operatorname{DF}(\mathcal{X},\mathcal{L})$ equals the weight of the induced $\mathbb{G}_m$-action on the fibre $\lambda_{\mathrm{CM}}(X,L)|_{t = 0}$ for $\lambda_{\mathrm{CM}}(X,L)$ the CM-line bundle that is determined by $(\mathcal{X},\mathcal{L})$.
\end{theorem}

\begin{theorem}\label{Ross:Thomas:Slope:Stab:DF}
Suppose that $L$ is a big and nef line bundle on an $n$-dimensional normal projective variety $X$.  Suppose that $Z \subsetneq X$ is a proper subscheme and with ideal sheaf $\Ish_Z$.  Let $\epsilon(L;Z)$ denote the nef threshold of $X$ with respect to $Z$.  For each $0 \leq t \leq \epsilon(L;Z)$ set
\begin{equation}\label{alpha0}
\alpha_0(L;Z; t) := \frac{((\pi^*L - t E)^n )}{n!}
\end{equation}
\begin{equation}\label{alpha1}
\alpha_1(L;Z; t) := - \frac{(\K_{X'} \cdot(\pi^*L - t E)^{n-1})}{2(n-1)!}
\end{equation}
and
\begin{equation}\label{slope}
\mu(X;L) := \frac{\alpha_1}{\alpha_0} = \frac{\alpha_1(0)}{\alpha_0(0)} \text{.}
\end{equation}
For each 
$c \in (0, \epsilon(L;Z)] \bigcap \QQ$
set
$$
\mu_c(L;\Ish_Z) := \frac{ \int_0^c \left( \alpha_1(L;Z;t) + \frac{\alpha_0'(L;Z;t)}{2} \right) \mathrm{d}t }{ \int_0^c \alpha_0(L;Z;t) \mathrm{d}t } \text{.}
$$
Fix a sufficiently divisible integer $m \in \ZZ$ so that $mc \in \ZZ$.  Then, with this notation and hypothesis the test configuration $(\mathcal{X}, \mathcal{L}_c^{\otimes m})$ has Donaldson-Futaki invariant $\operatorname{DF}(\mathcal{X},\mathcal{L}_c^{\otimes m})$ equal to some positive scalar multiple of $\mu(X;L) - \mu_c(L;\Ish_Z)$.
\end{theorem}

\begin{proof}
The idea is to use Corollary \ref{CM:line:bundle:epsilon:zero} together with the Riemann-Roch Theorem for normal varieties (e.g., in the form of \cite[Theorem A.1]{Boucksom-Hisamoto-Jonsson:2016}) to reduce the proof of the theorem to the classical case of Ross and Thomas \cite[Section 4]{Ross:Thomas:2007}.  

Fix a sufficiently small rational number $\epsilon > 0$ and consider the ample line bundle $L + \epsilon H$.  Suppose that 
$$c \in (0,\epsilon(L;Z)] \bigcap \QQ \text{.}$$  
Then $\pi^*L - c E$ is nef.  Also, 
$$\pi^*L \leq \pi^*(L+ \epsilon H)$$ 
whence 
$$0 \leq \pi^*L - c E \leq \pi^*(L+\epsilon H) - c E\text{.}$$  
Thus $\pi^*(L+\epsilon H) - c E$ is nef too.    In other words
$$
c \in (0,\epsilon(L+\epsilon H;Z)] \bigcap \QQ \text{.}
$$

Let $\mathcal{X} := \operatorname{DefNC}(X,Z)$ denote the deformation to the normal cone.   Replacing $m$ by some suitably large divisible integer $m$ if necessary, consider the test configurations 
$$(\mathcal{X}, p^*L^{\otimes m} - c m \mathcal{E})$$ 
and
$$(\mathcal{X}, p^*(L^{\otimes  m}  \otimes H^{\otimes \epsilon m}) - cm \mathcal{E}) ) \text{.}$$

Combining Corollary \ref{CM:line:bundle:epsilon:zero} and Theorem \ref{Paul:Tian:DF:Thm} it follows that
$$
\operatorname{DF}( (\mathcal{X}, p^*L^{\otimes m} - c m \mathcal{E}) ) = \lim\limits_{\epsilon \to 0} \operatorname{DF}( (\mathcal{X}, p^*(L^{\otimes m}  \otimes H^{\otimes \epsilon m}) - cm \mathcal{E})  ) 
$$

On other other hand, the conclusion of Theorem \ref{Ross:Thomas:Slope:Stab:DF} applied to $(\mathcal{X}, p^*(L^{\otimes  m}  \otimes H^{\otimes \epsilon m}) - cm \mathcal{E})  )$ follows as a consequence of the theory of Ross and Thomas \cite[Section 4]{Ross:Thomas:2007} combined  with the Riemann-Roch Theorem for normal varieties (e.g., in the form of \cite[Theorem A.1]{Boucksom-Hisamoto-Jonsson:2016}).

Finally, since
$$
\mu(X;L) - \mu_c(L;\Ish_Z) = \lim\limits_{\epsilon \to 0} (\mu(X;L+\epsilon H) - \mu_c(L+\epsilon H;\Ish_Z) )
$$
the conclusion desired by Theorem \ref{Ross:Thomas:Slope:Stab:DF} then follows.
\end{proof}

\begin{thebibliography}{99}
\bibitem[Arezzo]{Arezzo:DellaVedova:LaNave} C.~Arezzo, A.~Della~Vedova, and G.~La~Nave, \emph{Singularities and {$K$}-semistability}, Int. Math. Res. Not. IMRN (2012), no.~4, 849--869.
\bibitem[Boucks]{Boucksom-Hisamoto-Jonsson:2016} S.~Boucksom, T.~Hisamoto, and M.~Jonsson, \emph{Uniform {$\mathrm{K}$}-stability, {D}uistermaat-{H}eckman measures and singularities of pairs}, Ann. Inst. Fourier (Grenoble) \textbf{67} (2017), no.~2, 743--841.
\bibitem[Ross:T]{Ross:Thomas:2007} \bysame, \emph{A study of the {H}ilbert-{M}umford criterion for the stability of projective varieties}, J. Algebraic Geom. \textbf{16} (2007), no.~2, 201--255.
\end{thebibliography}
\end{document}
